\clearpage\fancyhead[L]{\small COMET.}\hypertarget{native-342-lower}{}\label{native-342-lower}
\begin{center}\textbf{COMET.}\par\vspace{4pt}\hrule\vspace{4pt}\fontsize{10}{12}\selectfont Göttingische gelehrte Anzeigen. Issue 151. pp. 1497\,..\,1499. 21 September 1811.\par\vspace{4pt}\hrule\vspace{4pt}\end{center}

\par\indent Our astronomers first saw this year’s great comet here on 22 August, well into the evening twilight. We had predicted its reappearance in issue 130 of our learned Notices [p. 327 of this volume]. Detailed observations were possible only later, since the city’s buildings obstruct the view north from the observatory too severely. Angular distances of the comet from fixed stars were measured. This method works very well for so bright a comet and gives a high degree of accuracy, especially when the observations are made with the instrument on a stand. As of 10 September, the only observations fully reduced are the following ones by Professor \textsc{Gauss}, made by this method.

\Needspace{9\baselineskip}
\begin{center}\fontsize{10}{12}\selectfont\setlength{\tabcolsep}{4pt}\renewcommand{\arraystretch}{1.15}\begin{tabular}{ll|c|c|c}
\multicolumn{2}{c|}{1811} & Mean time & right ascension & Northern declination \\\hline
Sept. & 4 & $8^{\mathrm u}\,28^{\mathrm m}\,47^{\mathrm s}$ & $158^\circ\;25'\;24''$ & $39^\circ\;18'\;2''$ \\
 & 6 & $8\;\;48\;\;38$ & $160\;\;\;23\;\;\;16$ & $40\;\;\;14\;\;\;16$ \\
 & 7 & $8\;\;57\;\;\;6$ & $161\;\;\;26\;\;\;10$ & $40\;\;\;41\;\;\;54$ \\
\end{tabular}
\end{center}

\par\indent Using these observations, Professor \textsc{Gauss} improved the provisional parabolic elements given at the place cited above and found the following results.

\Needspace{13\baselineskip}
\begin{center}\fontsize{10}{12}\selectfont\setlength{\tabcolsep}{4pt}\renewcommand{\arraystretch}{1.15}\begin{tabular}{l}
Perihelion passage 1811 Sept. 12.\\
\quad $5^{\mathrm u}\,21^{\mathrm m}\,15^{\mathrm s}$ mean time in Göttingen\\
\end{tabular}\end{center}
\begin{center}\fontsize{9}{11}\selectfont\setlength{\tabcolsep}{2pt}\renewcommand{\arraystretch}{1.15}\begin{tabular}{llr}
Longitude of perihelion & .\;.\;.\;.\;.\;.\;.\;.\;.\;.\;.\;.\;.\;.\;.\;. & $75^\circ\,17'\,34''$\\
Least distance from the Sun & .\;.\;.\;.\;.\;.\;.\;.\;.\;.\;.\;. & $1.04006$\\
Longitude of the ascending node & .\;.\;.\;.\;.\;.\;.\;.\;.\;.\;.\;. & $140^\circ\,24'\,13''$\\
Inclination of the orbit & .\;.\;.\;.\;.\;.\;.\;.\;.\;.\;.\;.\;.\;.\;.\;.\;. & $73\;7\;16$\\
\end{tabular}\end{center}
\noindent \emph{Retrograde} motion.

\par\indent So far there is no sign that the orbit is elliptical.

\par\indent To make the comet’s motion over the coming months easier to follow, Professor \textsc{Gauss} calculated a short ephemeris from the improved elements above. It shows that the comet comes closest to Earth on 15 October, when it is still nearly a quarter farther from Earth than the Sun is. If it is visible solely by\SourceLeaf{343} reflected light, it reaches its greatest brightness on 7 October. It is then about fifteen times as bright as it was on 11 April. On 3 October the comet passes less than one degree from the star $\eta$ in Ursa Major, and on 3 December less than one degree from the star $\alpha$ in Aquila.

\Needspace{15\baselineskip}\noindent Motion of the comet during the last months of this year

\begin{center}\fontsize{10}{12}\selectfont\setlength{\tabcolsep}{3pt}\renewcommand{\arraystretch}{1.15}\begin{tabular}{ll|c|c|c|c}
\multicolumn{2}{c|}{1811} & \shortstack{Right\\ascension} & \shortstack{Northern\\declination} & \shortstack{Distance\\from Earth} & \shortstack{Brightness} \\\hline
Sept. & 12\ $5^{\mathrm u}$ & $166^\circ\;54'$ & $42^\circ\;58'$ & 1.614 & 0.355 \\
 & 22\ $6$ & $182\;\;\;2$ & $47\;\;\;17$ & 1.443 & 0.432 \\
Oct. & 2\ $7$ & $203\;\;\;24$ & $49\;\;\;34$ & 1.306 & 0.490 \\
 & 12\ $8$ & $228\;\;\;45$ & $47\;\;\;16$ & 1.233 & 0.492 \\
\end{tabular}
\end{center}

\Needspace{15\baselineskip}\noindent Motion of the comet during the last months of this year

\begin{center}\fontsize{10}{12}\selectfont\setlength{\tabcolsep}{3pt}\renewcommand{\arraystretch}{1.15}\begin{tabular}{ll|c|c|c|c}
\multicolumn{2}{c|}{1811} & \shortstack{Right\\ascension} & \shortstack{Northern\\declination} & \shortstack{Distance\\from Earth} & \shortstack{Brightness} \\\hline
Oct. & 22\ $0^{\mathrm u}$ & $251^\circ\;23'$ & $39^\circ\;56'$ & 1.248 & 0.421 \\
Nov. & 1\ $10$ & $267\;\;\;58$ & $30\;\;\;32$ & 1.353 & 0.310 \\
 & 11\ $11$ & $279\;\;\;38$ & $22\;\;\;18$ & 1.531 & 0.208 \\
 & 21\ $12$ & $288\;\;\;13$ & $14\;\;\;55$ & 1.744 & 0.138 \\
Dec. & 1\ $13$ & $294\;\;\;54$ & $9\;\;\;42$ & 1.984 & 0.093 \\
 & 11\ $14$ & $300\;\;\;24$ & $5\;\;\;49$ & 2.232 & 0.064 \\
 & 21\ $15$ & $303\;\;\;22$ & $3\;\;\;0$ & 2.426 & 0.048 \\
\end{tabular}
\end{center}

{\rightskip=0pt plus 1em\par\indent The shape of the \emph{tail} is strange. It curves around the comet without being attached to it and ends in two branches inclined toward each other. There is, moreover, no trace at all of a \emph{nucleus} as such in the comet.\par}
